\section*{Capítulo \ref{ch:b2:cell-differentiation} --- Diferenciação celular: a célula muscular esquelética}

\begin{solution}{exo:b2:cell-differentiation:1}
\omterm{def:b2:cell-differentiation:differentiation}{Diferenciação}: a expressão estável dos genes e das funções de um tipo
celular. \omterm{def:b2:cell-differentiation:differentiation}{Determinação}: o comprometimento com um destino antes de ele
ser expresso. \omterm{def:b2:cell-differentiation:differentiation}{Potência}: o leque de destinos ainda abertos. Zigoto:
totipotente; célula da \omterm{def:b2:mammal-reproduction:implantation}{massa celular interna}: pluripotente; célula
satélite: unipotente (só músculo) --- ainda assim uma \omterm{def:b2:cell-differentiation:differentiation}{célula-tronco};
\omterm{def:b2:cell-differentiation:myogenesis}{fibra muscular}: diferenciada, pós-mitótica, sem \omterm{def:b2:cell-differentiation:differentiation}{potência}.
\end{solution}

\begin{solution}{exo:b2:cell-differentiation:2}
Célula do dermomiótomo (Pax3) $\to$ \omterm{def:b2:cell-differentiation:myogenesis}{mioblasto} induzido (\omterm{def:b2:cell-differentiation:myogenesis}{MyoD}, Myf5),
em divisão enquanto dura o FGF $\to$ saída do ciclo, miogenina $\to$
alinhamento e fusão num \omterm{def:b2:cell-differentiation:myogenesis}{miotubo} $\to$ genes musculares ativos,
\omterm{def:b2:cell-differentiation:fibre}{miofibrilas} montadas (MRF4) $\to$ fibra madura com núcleos periféricos,
mais células satélites quiescentes (Pax7).
\end{solution}

\begin{solution}{exo:b2:cell-differentiation:3}
Que o núcleo de uma célula diferenciada conserva todos os genes
necessários para formar um animal inteiro, e que o citoplasma do ovo
pode reiniciar seu programa: a \omterm{def:b2:cell-differentiation:differentiation}{diferenciação} é reversível e não é uma
perda de DNA. Ela não mostrou que todo núcleo diferenciado pode ser
reiniciado (a maioria das transferências fracassou), nem como a
reinicialização funciona, nem que a própria célula diferenciada
pudesse mudar de destino no lugar.
\end{solution}

\begin{solution}{exo:b2:cell-differentiation:4}
Discos Z delimitando o \omterm{def:b2:cell-differentiation:fibre}{sarcômero}; filamentos finos de actina ancorados
neles; filamentos grossos de miosina centrados na linha M; a banda A
(filamentos grossos) e a banda I (só filamentos finos). No
encurtamento, os filamentos finos deslizam em direção à linha M e os
discos Z se aproximam: a banda I e a zona H se estreitam, a banda A ---
o comprimento dos filamentos grossos --- não muda, e nenhum dos
filamentos encurta.
\end{solution}

\begin{solution}{exo:b2:cell-differentiation:5}
$0.2/2.5\times 10^{-6} = \num{80000}$ \omterm{def:b2:cell-differentiation:fibre}{sarcômeros}; $0.2/25\times 10^{-6}
= \num{8000}$ núcleos; um \omterm{def:b2:cell-differentiation:myogenesis}{mioblasto} por núcleo, portanto cerca de
\num{8000} \omterm{def:b2:cell-differentiation:myogenesis}{mioblastos} se fundiram (as células satélites acrescentam
mais depois).
\end{solution}

\begin{solution}{exo:b2:cell-differentiation:6}
O FGF induz a Id, que se liga à proteína parceira de que a \omterm{def:b2:cell-differentiation:myogenesis}{MyoD} precisa
para formar um dímero que se liga ao DNA; a \omterm{def:b2:cell-differentiation:myogenesis}{MyoD} está presente, mas
inativa, e as células continuam se dividindo. Quando o FGF é retirado,
a Id se degrada, a \omterm{def:b2:cell-differentiation:myogenesis}{MyoD} dimeriza, a miogenina é ativada, e as células
saem do ciclo e se fundem. \omterm{def:b2:cell-differentiation:myogenesis}{Mioblastos} que produzem Id de forma
constitutiva nunca se fundem: dividem-se indefinidamente e não formam
músculo.
\end{solution}

\begin{solution}{exo:b2:cell-differentiation:7}
Condição $0.4m = m^{2}/(1 + m^{2}) + 0.02$. Em $m = 2.1$: lado esquerdo
$0.84$, lado direito $4.41/5.41 + 0.02 = 0.835$: um estado
estacionário. Estado baixo: para $m$ pequeno, $0.4m \approx m^{2} + 0.02$,
de modo que $m \approx
0.055$ (esquerdo $0.022$, direito $0.023$). Elevada a 1, acima do
limiar perto de $0.41$, a célula sobe até $2.1$ e ali permanece: está
comprometida.
\end{solution}

\begin{solution}{exo:b2:cell-differentiation:8}
Com $k = 0.6$: em $m = 1$, a degradação $0.6$ supera a síntese
$0.52$; em $m = 0.5$, $0.30 > 0.22$; em $m = 2$, $1.2 > 0.82$: a reta
fica acima da curva em toda parte além do estado baixo, que é o único
estado estacionário. Uma célula cuja \omterm{def:b2:cell-differentiation:myogenesis}{MyoD} é degradada depressa demais
não consegue manter o estado ligado: ela recai no estado
indiferenciado, qualquer que seja a \omterm{def:b2:vertebrate-development:induction}{indução} --- nenhum músculo.
\end{solution}

\begin{solution}{exo:b2:cell-differentiation:9}
(a) Sem \omterm{def:b2:cell-differentiation:myogenesis}{mioblastos}, sem músculo esquelético: os dois fatores são
redundantes para a \omterm{def:b2:cell-differentiation:differentiation}{determinação}, de modo que ambos precisam faltar. (b)
Os \omterm{def:b2:cell-differentiation:myogenesis}{mioblastos} se formam e se alinham, mas nunca se fundem nem produzem
proteínas musculares: a miogenina é necessária para a \omterm{def:b2:cell-differentiation:differentiation}{diferenciação}.
(c) Quase normal, porque o Myf5 substitui a \omterm{def:b2:cell-differentiation:myogenesis}{MyoD} na \omterm{def:b2:cell-differentiation:differentiation}{determinação}.
\end{solution}

\begin{solution}{exo:b2:cell-differentiation:10}
Num fibroblasto, os genes musculares estão numa cromatina que a \omterm{def:b2:cell-differentiation:myogenesis}{MyoD}
consegue abrir; numa célula hepática, eles podem estar compactados em
heterocromatina, ou os próprios \omterm{prop:b2:cell-differentiation:transcriptional}{reguladores mestres} do fígado podem
reprimi-los, de modo que a célula não é competente. Teste: tratar
células hepáticas com um fármaco que afrouxa a cromatina (um inibidor
de histona-desacetilase) antes de acrescentar \omterm{def:b2:cell-differentiation:myogenesis}{MyoD}, ou acrescentar \omterm{def:b2:cell-differentiation:myogenesis}{MyoD}
com um fator que desloque o programa hepático, e medir a expressão de
miosina.
\end{solution}

\begin{solution}{exo:b2:cell-differentiation:11}
Disparos sustentados de baixa frequência elevam o cálcio intracelular
por longos períodos; vias ativadas pelo cálcio ligam os genes da
miosina lenta, da biogênese mitocondrial, da mioglobina e do
crescimento capilar, e desligam as isoformas rápidas. O programa mestre
da fibra (estado da \omterm{def:b2:cell-differentiation:myogenesis}{MyoD}, arquitetura do \omterm{def:b2:cell-differentiation:fibre}{sarcômero}) não é afetado: o
treino muda quais genes musculares são expressos, e não a identidade
da célula, e nenhum sinal vindo de um nervo consegue reabrir o
programa da cartilagem, fechado na \omterm{def:b2:cell-differentiation:differentiation}{determinação}.
\end{solution}

\begin{solution}{exo:b2:cell-differentiation:12}
O estado ligado de um fator autoativador é um ponto estável de uma
alça de retroalimentação, mantido por síntese contínua, e não por
alguma mudança nos genes; as marcas da cromatina acrescentam uma
segunda camada que mantém fechados os genes silenciados. A
reinicialização exige, portanto, levar o circuito abaixo de seu limiar
e reabrir a cromatina ao mesmo tempo, o que o citoplasma do ovo ou
quatro fatores só conseguem numa pequena fração das células: possível,
porque nada se perde, mas pouco eficiente, porque é preciso abrir duas
fechaduras independentes ao mesmo tempo.
\end{solution}

\begin{solution}{pb:b2:cell-differentiation:1}
\textbf{1.} $0.2/2.5\times 10^{-6} = \num{80000}$ \omterm{def:b2:cell-differentiation:fibre}{sarcômeros};
$0.2/25\times 10^{-6} = \num{8000}$ núcleos.
\textbf{2.} Cerca de \num{8000} por fibra; $4\times 10^{9}$ para o
músculo.
\textbf{3.} $\pi(30\times 10^{-6})^{2}\times 0.2 = \qty{5.7e-10}{m^3}$
por fibra; $2.8\times 10^{-4}$ \unit{m^3}, \qty{0.28}{L}, para o
músculo.
\textbf{4.} $5.7\times 10^{-10}/8000 = \qty{7e-14}{m^3} =
\qty{70000}{\micro m^3}$: o equivalente a 35 células comuns por núcleo.
\textbf{5.} O interior está repleto de \omterm{def:b2:cell-differentiation:fibre}{miofibrilas} que precisam correr
sem interrupção de uma ponta à outra; os núcleos na borda deixam a rede
contínua e ficam perto da membrana e de seus sinais.
\textbf{6.} $4\times 10^{9}/2000 = 2\times 10^{6}$: 21 duplicações,
10.5 dias.
\textbf{7.} $0.4m = m^{2}/(1 + m^{2}) + 0.02$. Em $m = 0.055$: lado
esquerdo $0.022$, lado direito $0.003 + 0.02 = 0.023$.
\textbf{8.} Em $m = 2.1$: esquerdo $0.84$, direito $0.815 + 0.02 = 0.835$.
\textbf{9.} Em $m = 0.41$: esquerdo $0.164$, direito $0.144 + 0.02 = 0.164$.
Logo abaixo, a síntese é menor que a degradação, de modo que $m$ cai
em direção a $0.055$; logo acima, a síntese supera a degradação e $m$
sobe em direção a $2.1$: qualquer deslocamento cresce --- instável.
\textbf{10.} $0.6 > 0.41$: a primeira célula vai para o estado ligado e
se diferencia; $0.3 < 0.41$: a segunda relaxa para $0.055$ e não se
diferencia.
\textbf{11.} Com $b = 0.4$, a curva de síntese começa em $0.4$ e fica
acima da reta $0.4m$ até $m \approx 3.3$: o cruzamento baixo desaparece.
A célula se diferencia espontaneamente, sem \omterm{def:b2:vertebrate-development:induction}{indução}.
\textbf{12.} A divisão reduz $m$ à metade; as filhas herdam $2.1/2 =
1.05$, acima do limiar $0.41$, e sobem de volta a $2.1$: o estado se
restabelece em cada filha. O comprometimento sobrevive enquanto o nível
reduzido à metade permanecer acima do limiar.
\textbf{13.} $4$ por mm $\times$ \qty{200}{mm} $= 800$ células satélites
por fibra; \qty{10}{\%} de \num{8000} núcleos, \num{800}, precisam ser
substituídos.
\textbf{14.} O segmento lesado contém $80$ células satélites; $800/80 =
10$, portanto 4 duplicações ($\times 16$, dando \num{1280}: 800 se
fundem, 480 restam para recompor o estoque) --- \qty{72}{h}, três dias.
\textbf{15.} Três dias de divisão contra duas a três semanas
observadas: o restante é a remoção dos detritos pelos macrófagos, o
atraso de ativação, a fusão, a montagem das \omterm{def:b2:cell-differentiation:fibre}{miofibrilas}, a reinervação e
a maturação do novo segmento.
\textbf{16.} A \omterm{def:b2:cell-differentiation:myogenesis}{MyoD} (e a miogenina) precisa religar os genes musculares
nas células ativadas; divisão e \omterm{def:b2:cell-differentiation:differentiation}{diferenciação} são excludentes --- a Id
precisa cair e o ciclo precisa parar antes que a miogenina possa agir
e a fusão ocorrer.
\textbf{17.} $0.95^{n} = 0.5$: $n = 13.5$, cerca de quatorze lesões.
\textbf{18.} $800\times 0.95^{20} = 800\times 0.36 \approx 290$.
\textbf{19.} Cada contração rompe fibras, cada ruptura recorre ao
estoque, e o estoque diminui alguns por cento a cada vez; após anos,
ele se esgota, as fibras rompidas não são reconstruídas e os
fibroblastos preenchem as lacunas com colágeno.
\textbf{20.} $(80/60)^{2} = 1.78$; núcleos $8000\times 1.78 = \num{14200}$:
cerca de \num{6200} a mais vindos das células satélites.
\textbf{21.} $0.28\times 1.78 = \qty{0.5}{L}$.
\textbf{22.} Mudaram: os genes das isoformas de miosina, a densidade de
mitocôndrias e de capilares, a mioglobina. Não mudaram: a identidade
muscular (estado da \omterm{def:b2:cell-differentiation:myogenesis}{MyoD}), a arquitetura do \omterm{def:b2:cell-differentiation:fibre}{sarcômero}, os núcleos e seu
genoma.
\textbf{23.} As fibras são pós-mitóticas e não conseguem se
multiplicar; os únicos caminhos são fibras mais grossas e mais núcleos
por fibra, vindos das células satélites.
\textbf{24.} Treino: espessamento limitado, pois a fibra não consegue
acrescentar núcleos; lesão: nenhum reparo --- os segmentos destruídos
são substituídos por tecido fibroso, como na distrofia.
\textbf{25.} \num{8000} núcleos por fibra; limiar $m \approx 0.41$;
cerca de três dias de divisão das células satélites para um reparo.
\end{solution}
